\subsection{序域的扩张}

定义 3.——设 K 是一个序域。K 的一个序扩张是一个二元组 (E, u)，其中 E 是一个序域，u 是从 K 到 E 的一个增同态。

设 K 为一个域，E 为一个序域，且设 \(u: K \rightarrow E\) 是一个同态。关系

\[
x \leqslant y \quad \text { if } \quad u (x) \leqslant u (y)
\]

是 K 上的一个全序关系，使 K 具有一个序域结构，称该结构由 E 的序诱导。若 K 与 E 都是序域，则同态 \(u: K \rightarrow E\) 是增同态，当且仅当 K 的序域结构由 E 的序域结构诱导。我们通常把 K 与它在 u 下的像等同起来。

例子. --- 1) 每个序域 K 都是 Q 的一个序扩张。事实上，K 是 Q 的扩张，因为其特征为零；另一方面，正如我们刚刚看到的，Q 只能以一种方式定序。

\begin{enumerate}
\def\labelenumi{\arabic{enumi})}
\setcounter{enumi}{1}
\tightlist
\item
  设 K 是一个序域，并设 K(X) 为 K 上关于单个未定元的有理函数域。让我们在多项式环 K{[}X{]} 上定义一个序：取正元素为 0 以及首项系数为正的那些多项式。这样我们得到一个全序环，其序延拓了 K 上的序。应用命题 2，我们赋予 K(X) 一个 K 的序扩张的结构。* 对 K = R，可以证明如此定义在 K(X) 上的序关系就是在 +∞ 附近的增长序（参见 VI，第 24 页，命题 4）。*
\end{enumerate}

定理 1. --- 对于 K 的扩张 E，要使 E 能够具有 K 的序扩张的结构，以下条件是必要且充分的：

（OE）关系 \(p_{1}x_{1}^{2} + \cdots + p_{n}x_{n}^{2} = 0\) 蕴含

\[
p _ {1} x _ {1} = \dots = p _ {n} x _ {n} = 0
\]

对元素对 \((x_{i},p_{i})\) 的任意有限序列，其中 \(x_{i}\) 属于 E，\(p_i\) 是 K 的正元素。

条件 (OE) 显然等价于：

(OE') 元素 \(-1\) 不是形如 \(px^2\) 的元素之和 ( \(x \in \mathbf{E}\) , \(p \in \mathbf{K}\) , \(p \geqslant 0\) )。

条件 (OE) 是必要的：如果 E 是 K 的一个序扩张，则元素 \(p_{i}x_{i}^{2}\) 在 E 中是正的，故当它们的和为零时它们为零。另一方面，\(p_{i}x_{i}^{2}=0\) 等价于 \(p_{i}x_{i}=0\) 。

反之，假设条件 (OE) 成立，则我们将通过在 E 中构造一个满足条件 \((\mathrm{AP}_{\mathrm{I}})\) , \((\mathrm{AP}_{\mathrm{II}})\) , \((\mathrm{AP}_{\mathrm{III}})\) 与 \((\mathrm{AP}_{\mathrm{IV}})\) 且包含 K 的正元素之集 \(K_{+}\) 的子集 P 来在 E 上定义一个序。这样的子集 P 必使 E 成为 K 的一个序扩张，因为我们会有 \(K \cap P = K_{+}\) ；事实上，如果 P 包含 K 的一个元素 -a \textless{} 0，则 a 将属于 \(P \cap (-P)\) ，这与 \((\mathrm{AP}_{\mathrm{III}})\)

相矛盾。为了定义 P，让我们考虑由 E 中满足 \((\mathrm{AP}_{\mathrm{I}})\) , \((\mathrm{AP}_{\mathrm{II}})\) 与 \((\mathrm{AP}_{\mathrm{III}})\) 且包含 \(K_{+}\) 与 E 的平方元素所成之集 C 的并的子集构成的集合 M。这个集合 M 是非空的，因为它包含集合 \(P_{0}\) ，即形如 \(\sum_{i} p_{i} x_{i}^{2}\) 的元素所成之集（\(P_{0}\) 满足 \((\mathrm{AP}_{\mathrm{III}})\) 这一事实由

（OE）立即可得）。

此外，M 是归纳的（集合论，III，第 154 页，定义 3）。因此，根据集合论，III，第 154 页的定理 2，M 中存在一个极大元，我们尚需证明它满足 \((\mathrm{AP}_{\mathrm{IV}})\) ；而这由以下引理得出：

引理。--- 设 \(P \in M\) 且 \(x \notin P\) ；则存在 \(P' \in M\) 使得 \(P \subset P'\) 且 \(-x \in P'\) 。

取 \(P' = P - xP\) ，并验证 \(P'\) 具有所需的性质。由于 \(0 \in C \subset P\) ，我们有 \(P \subset P'\) 。由此 \(C \subset P'\) 且 \(K_{+} \subset P'\) 。由于 \(1 \in C \subset P\) ，我们有 \(-x \in P'\) 。我们有

\[
\mathrm{P} ^ {\prime} + \mathrm{P} ^ {\prime} = \mathrm{P} - x \mathrm{P} + \mathrm{P} - x \mathrm{P} = \mathrm{P} + \mathrm{P} - x (\mathrm{P} + \mathrm{P}) \subset \mathrm{P} - x \mathrm{P} = \mathrm{P} ^ {\prime},
\]

由此得 \((\mathbf{AP}_1)\) 。我们有

\[
\begin{array}{r l} \mathrm {P^ {\prime} P^ {\prime}} = & (\mathrm{P} - x \mathrm{P}) (\mathrm{P} - x \mathrm{P}) \subset \\ & \subset \mathrm{PP} + x ^ {2} \mathrm{PP} - x (\mathrm{PP} + \mathrm{PP}) \subset \mathrm{P} + \mathrm{CP} - x \mathrm{P} \subset \mathrm{P} - x \mathrm{P} = \mathrm {P^ {\prime}}, \end{array}
\]

由此得 \((\mathrm{AP}_{\mathrm{II}})\) 。最后，验证 \((\mathrm{AP}_{\mathrm{III}})\) ：设给定一个形如 \(p - xq = -(r - xs)\) 的恒等式，其中 p、q、r、s 属于 P；由此我们推出关系 \(x(s + q) = p + r\) ；若 \(s + q \neq 0\) ，我们有

\[
x = (s + q) ^ {- 2} (s + q) (p + r) \in \mathrm{CPP} \subset \mathrm{P},
\]

与假设相反；因此 \(s + q = 0\) ，从而 \(p + r = 0\) ；由于 \(\mathbf{P}\) 满足 \((\mathbf{AP}_{\mathrm{III}})\) ，我们推断出 \(s = q = r = p = 0\) ，这就完成了证明。

推论 1（“阿廷–施赖尔定理”）。——在交换域 E 上存在一个序关系使其成为序域的充分必要条件是：关系 \(x_{1}^{2} + \cdots + x_{n}^{2} = 0\) 蕴含 \(x_{1} = \cdots = x_{n} = 0\) 。

必要性是显然的。反过来，所述条件蕴含 E 的特征为零，因此 E 是 Q 的扩张；于是条件 (OE) 得到满足，且定理 1 表明在 E 上存在 Q 的序扩张结构，即一个序域结构。

在 -1 为平方数的域 E 上不存在任何序域结构，特别地，在代数闭域上也是如此。

推论 2. --- 设 E 是 K 的一个扩张，且 E 容许具有 K 的序扩张的结构。元素 \(x \in E\) 在 E 上的每一个这样的结构下为正，当且仅当 x 具有形式 \(\sum_{i} p_{i} x_{i}^{2}\) ，其中 \(x_{i} \in E\) ，而

\(p_i\) 是 \(\mathbf{K}\) 的正元素。

该条件显然充分；它也是必要的，因为（采用定理 1 证明中的记号），如果 \(x \notin \mathbf{P}_0\) ，则存在一个极大元 \(\mathbf{P}\) 属于 \(\mathfrak{M}\) ，使得 \(x \notin \mathbf{P}\) ；于是根据引理，\(-x \in \mathbf{P}\) ，且 \(x\) 在由 \(\mathbf{P}\) 定义的序下不为正，因为 \(x \neq 0\) 。

